% !TeX root = ../../main.tex
\section{Introduction}
\label{sec:introduction}

Intersections of Hodge classes measure how the bundle of holomorphic
differentials varies over the moduli space of curves. Their
large-genus behavior depends on how the Chern indices vary with
the rank. The four results are grouped into two broad index
settings. One covers powers of the first class and fixed-index
insertions against that background; the other covers the top class
and indices at a fixed distance from it. The results give explicit
constants and
correction terms, together with the scaling that relates the top
class to each fixed defect.

Let $g,n$ be nonnegative integers with $2g-2+n>0$.  Denote by
$\MM_{g,n}$ the proper Deligne--Mumford stack over $\mathbb C$ of
stable curves of genus $g$ with $n$ marked points, and write $\MM_g$
when $n=0$ and $g\geq2$.  We work in its rational Chow ring
$A^*(\MM_{g,n};\mathbb Q)$, where degree means complex codimension.

The Hodge bundle $\EE$ has rank $g$ and fiber $H^0(C,\omega_C)$ at
$[C;p_1,\ldots,p_n]$, where $\omega_C$ is the dualizing sheaf of the
curve $C$.  Its Chern classes are the Hodge classes
\[\lambda_j=c_j(\EE)\in A^j(\MM_{g,n};\mathbb Q),
\qquad 1\leq j\leq g.\]
We also set $\lambda_0=1$ and $\lambda_j=0$ for every integer $j>g$.

The stack has complex dimension $3g-3+n$.  Integration selects the
component of that codimension and takes the degree of the resulting
zero-cycle, including stabilizer weights.  Scalar
coefficients may be extended when necessary.  For a positive integer
$\ell$, we study the \emph{pure Hodge integrals}
\begin{equation}
\int_{\MM_{g,n}}\lambda_{j_1}\cdots \lambda_{j_\ell},\quad j_1, \dots, j_\ell \in \{1, \dots, g\},
\label{eq:phi-1}
\end{equation}
as $g$ tends to infinity.

The integral in \eqref{eq:phi-1} vanishes unless
\[j_1+\cdots+j_\ell=3g-3+n.\]
For integers $0\leq m<n$ with $2g-2+m>0$, let
$\pi_{n,m}:\MM_{g,n}\to\MM_{g,m}$ forget the last $n-m$ markings
and stabilize.  The Hodge bundle on the source is pulled back from the
target.  When the degree condition above holds, the corresponding
product on the target has degree $3g-3+n>3g-3+m$, so it vanishes,
as does its pullback.  Hence \eqref{eq:phi-1} vanishes for $g\geq2$,
$n\geq1$ (take $m=0$), and for $g=1$, $n\geq2$ (take $m=1$).
In genus zero there are no positive-degree Hodge classes.  The only
remaining marked case is
\[\int_{\MM_{1,1}}\lambda_1=\frac{1}{24},\]
and the large-genus problem therefore reduces to the unmarked integrals
\begin{equation}
\int_{\MM_g}\lambda_{j_1}\cdots \lambda_{j_\ell},
\quad j_1, \dots, j_\ell \in \{1, \dots, g\},
\label{eq:phi-2}
\end{equation}
where $g\ge 2$.

For $g\geq2$ and an integer $h$ with $0\leq h\leq g$, set
\[a_g=\int_{\MM_g}\lambda_1^{3g-3},\quad b_g(h)=\int_{\MM_g}\lambda_{g-h} \lambda_1^{2g-3+h},
\quad b_g=b_g(0).\]
For a homogeneous polynomial $P(\lambda)$ of Hodge classes of degree
$d$, with $0\leq d\leq3g-3$, define
\[\langle P(\lambda) \rangle_g=\frac{1}{a_g}\int_{\MM_g} P(\lambda)\lambda_1^{3g-3-d}.\]

The geometric starting point is Mumford's
Grothendieck--Riemann--Roch description of the Hodge bundle
\cite[\S 5]{Mum83}, as recalled by
Faber--Pandharipande \cite[\S 1.1]{FP00}. It connects Hodge classes
to boundary geometry and permits two different approaches to
their asymptotics.

For Weil--Petersson volumes, Mirzakhani--Zograf established
fixed-marking expansions to every fixed order
\cite[Theorem 1.2]{MZ15}. We use the absolute leading normalization
proved by Liu--Xu \cite[Corollary 1.2]{LX26} and by Wang et al.\
\cite[Corollary 1.2]{W26}, together with the uniform vertex
estimates of Hide--Macera--Thomas \cite[Theorem 4.2]{HMT25}.
We also use Mirzakhani's comparisons, recalled in \cite[Lemma 14]{SW25},
and Shen--Wu's all-type bound \cite[Lemma 18]{SW25}.
Turning these inputs into pure Hodge asymptotics requires control
of the entire stable-graph sum: arbitrarily many self-loops
contribute at leading order.

The first two results determine the absolute asymptotics of $a_g$
and the effect of any fixed list of Hodge insertions.

\begin{thm}
\label{thm:lambda-one}
For real $x\geq0$, define
\[\Cin(x)=\int_0^x \frac{1-\cos t}{t}\dd t, \qquad
N_g=\frac{(2g-3)!(3g-3)!(4\pi^2)^{2g-3}}{\sqrt{g}(24\pi^2)^{3g-3}}.\]
As $g\to \infty$, we have
\begin{equation}
a_g=\frac{e^{\Cin(2\pi)}}{\sqrt{\pi}}N_g
\left(1+\frac{9}{8}g^{-1}+\left(\frac{257}{128}+\frac{1}{8\pi^2}\right)g^{-2}+O(g^{-3})\right).
\label{eq:111}
\end{equation}
\end{thm}

\begin{thm}
\label{thm:fixed-insertions}
Fix $\ell$ and the positive indices $j_1,\ldots,j_\ell$, and put
$P(\lambda)=\lambda_{j_1}\cdots\lambda_{j_\ell}$.  For this list of
indices, define
\[F(j)=\prod_{k=1}^\ell j_k!, \quad G(j)=\sum_{k=1}^\ell j_k(j_k-1)(j_k-2).\]
Then, as $g\to\infty$, we have
\begin{equation}
\langle P(\lambda) \rangle_g=\frac{1}{F(j)}\left(1+\frac{4\pi^2-42}{9g^2}G(j)+
\frac{4\pi^2-40}{3g^3}G(j)+O(g^{-4})\right).
\label{eq:P111}
\end{equation}
\end{thm}

The next two results concern the opposite range of indices.
The inserted class is $\lambda_g$, or $\lambda_{g-h}$ with $h$
fixed while $g$ increases. Their scale is governed by the first
positive zero of a Bessel function.

\begin{thm}
\label{thm:top-hodge}
Let $g_2=2g-3$, $J_n(z)$ be the $n$-th Bessel function of the first kind, $z_{0,1}$ be the first
positive zero of $J_0(z)$, and $\rho=6z_{0,1}J_1(z_{0,1})$.
Set
\[b=-\frac{1}{4}+\frac{1}{6z_{0,1}^2}+\frac{\pi^2}{6}J_1(z_{0,1})^2.\]
As $g\to\infty$, we have
\begin{equation}
b_g=\frac{1}{J_1(z_{0,1})}\frac{2}{(2\pi)^{2g}}\frac{(g_2!)^2}{\rho^{g_2}\sqrt{2\pi g_2}}\left(
1+\frac{b}{g_2}+O(g_2^{-2})\right).
\label{eq:g111}
\end{equation}
\end{thm}

\begin{thm}
\label{thm:fixed-defect}
With $\rho$ as above, for a fixed integer $h\geq0$, we have
\begin{equation}
\frac{b_g(h)}{b_g}\sim \frac{1}{h!}\left(\frac{24g^2}{\rho^2}\right)^h,\quad g\to\infty.
\label{eq:gh11}
\end{equation}
\end{thm}

Theorem~\ref{thm:fixed-insertions} shows that the normalized product
approaches the reciprocal factorial product, with no term of order
$g^{-1}$. The first correction depends only on the cubic sum
$G(j)$. In Theorem~\ref{thm:fixed-defect}, each fixed increase of
the defect contributes two powers of $g$; the factor $1/h!$
comes from ordered transport integrals. All insertion indices and
defects in these statements are fixed. The last theorem asserts
asymptotic equivalence, without a relative error rate.

Sections~3--5 prove Theorems~\ref{thm:lambda-one}
and~\ref{thm:fixed-insertions} by a stable-graph expansion.
The boundary identity of Andersen et al.\
\cite[Lemma 3.10]{ABCDGLW23} converts the geometric formulas into
products of Weil--Petersson vertex integrals. Section~3 first
resums the self-loops, proves an absolute tail bound for all graphs,
and reduces each fixed asymptotic order to finitely many cores
with summable remainders. Section~4 evaluates the terms needed
for the second correction in Theorem~\ref{thm:lambda-one}.
Section~5 introduces labeled sources for the insertions and
controls the resulting decorated graph sums.

Appendix~A supplies the finite calculations used in this route.
They include the descendant coefficients of Wang et al.\
\cite[Lemma 7.2 and Appendix C.1]{W26}, correlated loop moments
and the core contributions. The required second absolute volume
coefficient is reconstructed from the exact recursions of
Mirzakhani--Zograf \cite[\S 2]{MZ15}. Its agreement with the
conditional resurgent prediction of Eynard et al.\
\cite[\S 6.1, (6.10)]{EGGLS24} is an a posteriori check;
the reconstruction itself is unconditional.

Theorem~\ref{thm:top-hodge} has an independent proof in Section~6.
The $\lambda_g$ formula of Faber--Pandharipande
\cite[Theorem 1]{FP03} gives exact vertex integrals, and the
restriction of the top class eliminates every graph with a cycle.
The genus-zero inverse relation recalled by Kaufmann--Manin--Zagier
\cite[(0.8)]{KMZ96v1} and the Bessel analysis of Manin--Zograf
\cite[\S 6, (17) and Proof 6.3]{MaZ00} provide the analytic
background. Uniform bounds for all stable vertex types control
the whole tree sum. A saddle calculation and the elliptic-leaf
contribution then give the displayed correction.

For Theorem~\ref{thm:fixed-defect}, Sections~7--8 use Buryak's relation
between the linear Hodge potential and the ILW equation
\cite[Theorem 1.1]{Bur15}. Section~7 derives exact transport and
Borel identities, including the free-energy integration constants.
Section~8 identifies the principal singularity at each defect and
evaluates its transport constant. The common complex domains and
uniform kernel and primitive estimates are proved in Appendix~B.
Both dominant singularities enter the coefficient transfer,
following the method of Flajolet--Odlyzko \cite[\S 2]{FO90}.
A nef intersection argument controls the complete change of
background needed to return to $b_g(h)$.

Section~2 collects the geometric identities used by both routes.
After that section, readers interested primarily in
Theorems~\ref{thm:top-hodge} and~\ref{thm:fixed-defect}
may proceed directly to Sections~6--8; these sections do not use
the Weil--Petersson correction coefficients. Appendix~C records
exact normalization checks in genera two and three, together with
reproducibility instructions for the finite calculations. Such checks
support the analytic proofs without replacing their uniform
remainder estimates.
